\subsection{DIRECT LIMIT OF TENSOR ALGEBRAS}

Let \((\mathbf{A}_{\alpha},\phi_{\beta \alpha})\) be a directed direct system of commutative rings and \((\mathbf{M}_{\alpha},f_{\beta \alpha})\) a direct system of \(\mathbf{A}_{\alpha}\) -modules (II, § 6, no. 2); let \(\mathbf{A} = \varinjlim \mathbf{A}_{\alpha}\) and \(\mathbf{M} = \varinjlim \mathbf{M}_{\alpha}\) , which is an A-module. For \(\alpha \leqslant \beta\) an \(\mathbf{A}_{\alpha}\) -algebra homomorphism (no. 3, formula (8)) \(f_{\beta \alpha}' = \mathsf{T}(f_{\beta \alpha})\colon \mathsf{T}_{\mathbf{A}_{\alpha}}(\mathbf{M}_{\alpha})\to \mathsf{T}_{\mathbf{A}_{\beta}}(\mathbf{M}_{\beta})\) is derived canonically from the \(\mathbf{A}_{\alpha}\) -homomorphism \(f_{\beta \alpha}\colon \mathbf{M}_{\alpha}\to \mathbf{M}_{\beta}\) and it follows from (9) (no. 3) that \((\mathsf{T}_{\mathbf{A}_{\alpha}}(\mathbf{M}_{\alpha}),f_{\beta \alpha}')\) is a direct system of \(\mathbf{A}_{\alpha}\) -algebras. On the other hand let \(f_{\alpha}\colon \mathbf{M}_{\alpha}\to \mathbf{M}\) be the canonical \(\mathbf{A}_{\alpha}\) -homomorphism; an \(\mathbf{A}_{\alpha}\) -algebra homomorphism \(f_{\alpha}'\colon \mathsf{T}_{\mathbf{A}_{\alpha}}(\mathbf{M}_{\alpha})\to \mathsf{T}_{\mathbf{A}}(\mathbf{M})\) is derived (no. 3, formula (8)) and it also follows from (9) (no. 3) that the \(f_{\alpha}'\) constitute a direct system of \(\mathbf{A}_{\alpha}\) -homomorphisms.

PROPOSITION 6. The A-homomorphism \(f' = \varinjlim f_{\alpha}' : \varinjlim T_{A_{\alpha}}(M_{\alpha}) \to T_{A}(M)\) is a graded algebra isomorphism.

To simplify we write \(\mathbf{E} = T_{\mathbf{A}}(\mathbf{M})\) and \(\mathbf{E}' = \varinjlim T_{\mathbf{A}_{\alpha}}(\mathbf{M}_{\alpha})\) and let

\[
g _ {\alpha}: \mathsf {T} _ {\mathrm{A} _ {\alpha}} (\mathbf {M} _ {\alpha}) \rightarrow \mathbf {E} ^ {\prime}
\]

be the canonical \(\mathbf{A}_{\alpha}\) -homomorphism. Clearly the composite \(\mathbf{A}_{\alpha}\) -linear mappings \(\mathbf{M}_{\alpha} \xrightarrow{\phi_{\mathbf{M}_{\alpha}}} \mathsf{T}_{\mathbf{A}_{\alpha}}(\mathbf{M}_{\alpha}) \xrightarrow{g_{\alpha}} \mathbf{E}'\) form a direct system and there is therefore one and only one A-linear mapping \(u = \varinjlim (g_{\alpha} \circ \phi_{\mathbf{M}_{\alpha}}): \mathbf{M} \to \mathbf{E}'\) such that

\[
u \circ f _ {\alpha} = g _ {\alpha} \circ \phi_ {M _ {\alpha}}
\]

for all \(\alpha\) . This mapping itself factorizes uniquely (no. 1, Proposition 1) into \(\mathbf{M} \xrightarrow{\phi_{\mathbf{M}}} \mathrm{E} \to \mathrm{E}'\) , where \(h\) is an A-algebra homomorphism. It will suffice to prove that \(h \circ f' = 1_{\mathrm{E}'}\) and \(f' \circ h = 1_{\mathrm{E}}\) .

To this end note that, for every index \(\alpha\) , (no. 3, formula (8))

\[
h \circ f _ {\alpha} ^ {\prime} \circ \phi_ {M _ {\alpha}} = h \circ \phi_ {M} \circ f _ {\alpha} = u \circ f _ {\alpha} = g _ {\alpha} \circ \phi_ {M _ {\alpha}}
\]

whence, by the uniqueness assertion of no. 1, Proposition 1,

\[
h \circ f _ {\alpha} ^ {\prime} = g _ {\alpha}
\]

for all \(\alpha\) ; it follows that \((h \circ f') \circ g_{\alpha} = g_{\alpha}\) for all \(\alpha\) and hence \(h \circ f' = 1_{\mathbf{E}'}\) by definition of a direct limit.

On the other hand, by virtue of no. 3, formula (8),

\[
f ^ {\prime} \circ u \circ f _ {\alpha} = f ^ {\prime} \circ g _ {\alpha} \circ \phi_ {M _ {\alpha}} = f _ {\alpha} ^ {\prime} \circ \phi_ {M _ {\alpha}} = \phi_ {M} \circ f _ {\alpha},
\]

whence again \(f' \circ u = \phi_{\mathbb{M}}\) by definition of a direct limit; we conclude that \(f' \circ h \circ \phi_{\mathbb{M}} = \phi_{\mathbb{M}}\) and the uniqueness property of no. 1, Proposition 1 gives \(f' \circ h = 1_{\mathbb{E}}\) .

Proposition 6 can also be shown by observing that, for every integer \(n \geqslant 1\) , there is a canonical A-module isomorphism \(\lim_{\longrightarrow} \mathsf{T}_{\mathbf{A}_{\alpha}}^{n}(\mathbf{M}_{\alpha}) \to \mathsf{T}_{\mathbf{A}}^{n}(\mathbf{M})\) , as follows by induction on \(n\) from II, § 6, no. 3, Proposition 7. It is immediately verified that these isomorphisms are the restrictions of \(f'\) .

